\documentclass[11pt]{article}
\usepackage{amsmath,amssymb,amsfonts}
\usepackage{physics}
\usepackage{geometry}
\usepackage{hyperref}
\usepackage{cite}
\usepackage{bm}
\title{Coordinate Eigenstates as Distributional Sections}
\begin{document}
\maketitle


\section{Geometric Kinematics}

\begin{axiom}[Configuration manifold]
The configuration space of a nonrelativistic particle is a smooth manifold
$\Omega$ (possibly with boundary), equipped with its Borel $\sigma$-algebra.
\end{axiom}

\begin{axiom}[State bundle]
Associated to $\Omega$ is a complex Hermitian line bundle
\[
\pi : \mathcal L \to \Omega ,
\]
called the \emph{state bundle}.
\end{axiom}

\begin{axiom}[Physical states]
Physical states are square-integrable sections
\[
\psi \in \Gamma_{L^2}(\mathcal L),
\]
defined up to no equivalence relation (states are vectors, not rays).
\end{axiom}

\begin{remark}
The usual Hilbert space $L^2(\Omega)$ arises only after trivializing
$\mathcal L$. All geometric statements below are trivialization-independent.
\end{remark}

---

\section{Position as a Base-Space Observable}

\begin{axiom}[Position observable]
The position observable is the canonical projection
\[
X : \mathcal L \to \Omega ,
\]
acting on sections by pullback along evaluation.
\end{axiom}

\begin{definition}[Position POVM]
For any measurable $\Delta \subset \Omega$, define
\[
(E_X(\Delta)\psi)(x) = \chi_\Delta(x)\psi(x),
\]
where $\chi_\Delta$ is the indicator function on the base manifold.
\end{definition}

\begin{proposition}
$E_X$ is a projection-valued measure depending only on the base-space structure
of $\mathcal L$.
\end{proposition}

---

\section{Distributional Sections}

\begin{definition}[Test section space]
Let $\Phi \subset \Gamma_{L^2}(\mathcal L)$ be a dense space of smooth,
compactly supported sections.
\end{definition}

\begin{definition}[Distributional section]
A distributional section is a continuous linear functional
\[
T : \Phi \to \mathbb C.
\]
The space of such sections is denoted $\Gamma'(\mathcal L)$.
\end{definition}

\begin{remark}
$\Gamma'(\mathcal L)$ is the geometric analogue of the dual space
$\Phi^\times$ in rigged Hilbert space theory.
\end{remark}

---

\section{Coordinate Eigenstates as Delta Sections}

\begin{definition}[Coordinate eigensection]
For any point $x_0 \in \Omega$, the coordinate eigenstate is the
distributional section $\delta_{x_0} \in \Gamma'(\mathcal L)$ defined by
\[
\langle \delta_{x_0}, \psi \rangle := \psi(x_0).
\]
\end{definition}

\begin{theorem}[Eigenvalue property]
For all smooth functions $f$ on $\Omega$,
\[
f(X)\,\delta_{x_0} = f(x_0)\,\delta_{x_0}
\]
in the distributional sense.
\end{theorem}

\begin{corollary}
Coordinate eigenstates are eigenvectors of all position observables
but belong to $\Gamma'(\mathcal L)\setminus\Gamma_{L^2}(\mathcal L)$.
\end{corollary}

---

\section{Local Trivializations and Dirac Notation}

In any local trivialization of $\mathcal L$, one may write
\[
\langle x | x_0 \rangle = \delta(x-x_0),
\]
recovering the familiar Dirac formalism introduced by
\emph{:contentReference[oaicite:0]{index=0}}.

\begin{remark}
The delta function is not a function on $\Omega$ but a section-valued
distribution supported on a single base point.
\end{remark}

---

\section{Resolution of the Identity}

\begin{theorem}[Geometric completeness]
For any $\psi \in \Gamma_{L^2}(\mathcal L)$,
\[
\psi = \int_\Omega \delta_x \, \psi(x)\, dx,
\]
where the integral is understood weakly as a reconstruction of sections.
\end{theorem}

\begin{corollary}
The family $\{\delta_x\}_{x\in\Omega}$ provides a continuous geometric frame
for $\Gamma_{L^2}(\mathcal L)$.
\end{corollary}

---

\section{Bounded Configuration Spaces}

\begin{theorem}[Bounded base manifolds]
Let $\Omega$ be compact or bounded with boundary.
Then:
\begin{enumerate}
\item The state bundle $\mathcal L$ remains well-defined.
\item Delta sections $\delta_x$ exist for all $x\in\Omega$.
\item The position observable remains sharply defined.
\item No boundary conditions are required for position.
\end{enumerate}
\end{theorem}

\begin{remark}
This sharply contrasts with momentum observables, which require additional
bundle connections and boundary conditions and may fail to exist globally.
\end{remark}

---

\section{Momentum as a Connection-Dependent Observable}

\begin{axiom}[Connection]
Momentum observables arise from a Hermitian connection
\[
\nabla : \Gamma(\mathcal L) \to \Gamma(T^*\Omega \otimes \mathcal L).
\]
\end{axiom}

\begin{proposition}
Momentum eigenstates exist only when $\nabla$ admits global parallel sections.
\end{proposition}

\begin{corollary}
On bounded $\Omega$, momentum eigenstates generally do not exist as
distributional sections.
\end{corollary}

---

\section{Measurement and Collapse}

\begin{axiom}[Operational localization]
A finite-resolution position measurement corresponds to restricting a section
to a measurable region $\Delta\subset\Omega$.
\end{axiom}

\begin{theorem}[No geometric collapse]
No physical process produces a delta section $\delta_x$.
Exact localization corresponds to a singular limit in $\Gamma'(\mathcal L)$.
\end{theorem}

\begin{remark}
Collapse is not a geometric flow on $\mathcal L$ but an update of information
conditioned on a base-space event.
\end{remark}

---

\section{Interpretive Summary}

\begin{itemize}
\item Coordinate eigenstates are delta-supported sections of the state bundle.
\item They are geometrically well-defined but physically unrealizable.
\item Position is a base-manifold observable, not a generator of motion.
\item The bundle formalism makes the asymmetry between position and momentum
structural rather than interpretive.
\end{itemize}
\end{document}
